\documentclass[11pt]{article}

\usepackage[margin=1in]{geometry}
\usepackage{amsmath}
\usepackage{amssymb}
\usepackage{bm}
\usepackage{dsfont}
\usepackage[T1]{fontenc}
\usepackage{textcomp}
\usepackage{hyperref}

% --- review annotations -------------------------------------------------
% Comment out the next line (and uncomment the one after) for a clean build.
\usepackage[textwidth=3.2cm]{todonotes}
%\usepackage[disable]{todonotes}
\newcommand{\jc}[1]{\todo[inline,color=yellow!40]{#1}}
\newcommand{\jm}[1]{\todo[color=orange!40]{#1}}
% ------------------------------------------------------------------------

\newcommand{\rr}{\mathbf{r}}
\newcommand{\ee}{\hat{\mathbf{e}}}
\newcommand{\KK}{\mathbf{K}}
\newcommand{\LL}{\Lambda}
\newcommand{\re}{\rr^{e}}

\title{An elastic network model with $k_{ij}=k(l_{ij})$}
\author{Julian Echave}
\date{2026-08-26}

\begin{document}

\maketitle

\begin{abstract}
This note sets up an elastic network model of proteins whose goal is to
predict how mutations change structure, motions, and stability. In standard
ENMs the parameters are fixed by a known structure: each spring's rest length
is the distance between its residues in that structure, so modeling a mutant
requires supplying the mutant structure, which is itself one of the things to
be predicted. Here, instead, we tie spring stiffnesses to rest lengths,
$k(l_{ij})$, so the rest lengths alone define the model and the structure
follows by minimization. A mutation is then just a perturbation of the rest
lengths at a site, and the mutant's structure, motions, and thermodynamic
changes (energy, entropy, and free energy) are derived, not assumed. A side
effect is that networks can now carry strain and frustration, which real
proteins are known to display but standard ENMs exclude by construction. We
work out the full formalism and end with a list of open questions.
\end{abstract}

% \tableofcontents

\section{Introduction}
\label{sec:intro}

A protein can be modelled as an elastic network of nodes, representing amino
acids, connected by springs, representing interactions.
Such an elastic network model (ENM) is completely specified by the spring force
constants $k_{ij}$ and the spring equilibrium lengths $l_{ij}$.
Typically, given a native conformation $\re$, the ENM is set up so that
$l_{ij} = d_{ij}(\re)$ and $k_{ij} = k(d_{ij}(\re))$.
The structure thus enters twice: it fixes the parameters, and it is the
conformation at which the model is relaxed.

We would like to use such a model to predict how a mutation changes a protein:
its structure, its motions, its stability.
But a mutation changes the parameters, and here the parameters are the structure.
To mutate the model one must hand it the mutant structure, which is what we
wanted it to predict.
The same identification has a second consequence: every spring is relaxed at
$\re$, so the network carries no strain, and frustration, which proteins are
known to display, cannot be represented at all.

Both problems come from tying the parameters to a structure, so we untie them.
We take the force constants to be a function of the equilibrium lengths of the
springs, $k_{ij} = k(l_{ij})$, rather than of the distances between nodes in a
given structure.
The model is then determined by the $l_{ij}$ alone.
The springs need not be relaxed at the equilibrium conformation, but may be
strained, so the network may be frustrated; and a mutation of a site is modelled
by perturbing the equilibrium lengths of the springs connected to it, with the
mutant structure following from the model rather than being supplied to it.

This note derives the potential energy function, its minimum, the conformational
distribution, and thermodynamic properties, first for a general network and
then for a reference and its mutant, ending with the mutational changes.

\section{The model}

\subsection{Nodes}

A protein is represented by $N$ nodes, one per residue, joined in pairs by springs.
A conformation is the set of the node positions,
\begin{equation}
\rr = (\rr_1, \rr_2, \ldots, \rr_N) \in \mathbb{R}^{3N}.
\end{equation}
We write
\begin{equation}
\rr_{ij} \;\equiv\; \rr_j - \rr_i
\end{equation}
for the vector that goes from node $i$ to node $j$.
The distance between the two nodes in conformation $\rr$ is the norm of this vector:
\begin{equation}
d_{ij}(\rr) \;=\; \lVert \rr_{ij} \rVert.
\label{eq:dij}
\end{equation}

\subsection{Springs}

Every pair of nodes is joined by a spring.
A spring properties are its force constant $k_{ij}$ and its equilibrium length $l_{ij}$.
We assume the force constant of a spring to be determined by its own equilibrium length through a function $k$,
\begin{equation}
k_{ij} \;=\; k(l_{ij}),
\qquad k : [0,\infty) \to [0,\infty).
\label{eq:kl}
\end{equation}
We assume of $k$ only that it is non-increasing, so that long springs are weakly coupled or uncoupled, and that $k(l) = 0$ for $l$ beyond some range, so that springs long enough to represent no contact do not interact. \jm{Do we really assume this?}

Since the $k_{ij}$ are not independent parameters, the model is determined by the set of length parameters:
\begin{equation}
\LL \;=\; \{\, l_{ij} \,:\, 1 \le i < j \le N \,\},
\qquad
|\LL| = \tfrac12 N(N-1).
\end{equation}


This is the one departure from standard practice, and it is worth stating
carefully, because in the standard case the departure is invisible. A standard
elastic network model is built on a given structure, which it takes as the
equilibrium conformation $\re$, and assigns each spring the equilibrium length
$l_{ij} = d_{ij}(\re)$ it already has in that structure. Its force constants are
then written as a function of $d_{ij}(\re)$ --- but since $l_{ij}$ and
$d_{ij}(\re)$ are equal there, a function of one is a function of the other, and
nothing distinguishes $k(l_{ij})$ from $k(d_{ij}(\re))$. Here the $l_{ij}$ are left
free rather than read off a structure, the two cease to coincide, and
\eqref{eq:kl} states which of them $k$ is a function of: the parameter $l_{ij}$.


\subsection{The potential}

The potential energy of a conformation $\rr$ for a network with parameters $\LL$
is a sum of harmonic terms, one per spring,
\begin{equation}
\boxed{\;
V(\rr; \LL) \;=\; \tfrac12 \sum_{i<j} k(l_{ij})\,\big(d_{ij}(\rr) - l_{ij}\big)^{2}.
\;}
\label{eq:V}
\end{equation}
The sum runs over all $\tfrac12 N(N-1)$ springs, and $V$ evaluated at any
conformation is the elastic energy stored in them. Every term is non-negative, so
$V \ge 0$, and $V(\rr;\LL) = 0$ requires every spring to sit at its own equilibrium
length, $d_{ij}(\rr) = l_{ij}$ for every active spring at once. Whether any
conformation achieves this is a question about $\LL$, not a property of the model:
for some parameter sets no conformation does, and then $V > 0$ everywhere. Section
\ref{sec:frustration} takes this up.

A spring with $k(l_{ij}) > 0$ contributes to $V$ and is said to be \emph{active}; a
spring with $k(l_{ij}) = 0$ contributes nothing. The set of active springs is the
network's contact map. Since $k$ is a function of $l_{ij}$, which springs are
active is fixed by $\LL$ alone: the contact map is a property of the parameters,
and does not change as the network moves. In a standard model, where
$l_{ij} = d_{ij}(\re)$, this same set would be read off the structure instead.

\section{The minimum of the potential}

\subsection{The equilibrium condition}

To find the conformations that minimise \eqref{eq:V} we need the gradient of $V$,
and hence the gradient of the distances. Let
\begin{equation}
\ee_{ij} \;\equiv\; \frac{\rr_{ij}}{d_{ij}}
\label{eq:eij}
\end{equation}
be the unit vector along the spring, pointing from $i$ to $j$; it satisfies
$\ee_{ij} = -\ee_{ji}$. Differentiating $d_{ij}^2 = \rr_{ij}\cdot\rr_{ij}$ with
respect to $\rr_i$, and using $\partial \rr_{ij}/\partial\rr_i = -\mathbf{I}$,
\begin{equation}
2\, d_{ij} \frac{\partial d_{ij}}{\partial \rr_i} \;=\; -\,2\,\rr_{ij},
\qquad\text{so}\qquad
\frac{\partial d_{ij}}{\partial \rr_i} \;=\; -\frac{\rr_{ij}}{d_{ij}}
\;=\; -\,\ee_{ij}.
\label{eq:ddij}
\end{equation}
The derivative with respect to $\rr_j$ follows the same way, now with
$\partial \rr_{ij}/\partial\rr_j = +\mathbf{I}$, giving
$\partial d_{ij}/\partial \rr_j = +\ee_{ij}$. Moving node $i$ away from node $j$
lengthens the spring, and \eqref{eq:ddij} records that: the distance decreases
fastest along $-\ee_{ij}$, which points from $j$ back towards $i$.

Differentiating \eqref{eq:V} by the chain rule, and noting that $k(l_{ij})$ is a
constant with respect to $\rr$,
\begin{equation}
\frac{\partial V}{\partial \rr_i}
\;=\; -\sum_{j \ne i} k(l_{ij})\,\big(d_{ij}(\rr) - l_{ij}\big)\,\ee_{ij}.
\label{eq:grad}
\end{equation}
Each term lies along the line joining the two nodes, so the force it exerts is
central: it has no component transverse to the spring. That $k$ carries no
derivative is a consequence of \eqref{eq:kl}. Had the force constants been written
as a function of the equilibrium distance, $k_{ij} = f(d_{ij}(\re))$, they would
still be constants with respect to $\rr$ and this gradient would be unchanged; what
\eqref{eq:kl} buys is not simplicity here but the freedom, used in
Section~\ref{sec:mutation}, to vary $l_{ij}$ and have $k_{ij}$ follow.

Given the parameters $\LL$, an equilibrium conformation $\re$ is one at which the
gradient vanishes,
\begin{equation}
\boxed{\;
\sum_{j \ne i} k(l_{ij})\,\big(d_{ij}(\re) - l_{ij}\big)\,\ee^{\,e}_{ij} \;=\; \mathbf{0},
\qquad i = 1,\ldots,N,
\;}
\label{eq:equilibrium}
\end{equation}
where the overall sign of \eqref{eq:grad} has been dropped, the right-hand side
being zero. These are $3N$ scalar equations for the $3N$ unknown coordinates, of
which $3N-6$ are independent: $V$ depends on the conformation only through the
distances \eqref{eq:dij}, so it is invariant under the three translations and three
rotations, and along those six directions the gradient vanishes identically.

One consequence of that invariance should be recorded now, because it matters
later: $\re$ is determined only up to a rigid motion.

Separately, and for a different reason, $V$ is also invariant under reflection.
This does not follow from invariance under rigid motions --- a reflection is not
one --- but from the same fact as before, that $V$ sees the conformation only
through the distances \eqref{eq:dij}, and reflecting a conformation leaves every
$d_{ij}$ unchanged. A chiral $\re$ therefore has a mirror image of equal energy,
and the minimum of $V$ is not unique. Below, $\re$ denotes one chosen minimum, and
every quantity derived from it is understood relative to that choice.

\subsection{Strain}

The combination $d_{ij}(\re) - l_{ij}$ appearing in \eqref{eq:equilibrium} is the
amount by which a spring departs from its equilibrium length at the equilibrium
conformation. It is convenient to name it. Let $\re$ be an equilibrium
conformation and define, for every spring,
\begin{equation}
s_{ij} \;\equiv\; d_{ij}(\re) - l_{ij},
\label{eq:sij}
\end{equation}
its \emph{strain}. A spring with $s_{ij} > 0$ is longer than its equilibrium length
and so is stretched, one with $s_{ij} < 0$ is compressed, and one with $s_{ij} = 0$
is relaxed.

In these terms the equilibrium condition \eqref{eq:equilibrium} reads
\begin{equation}
\sum_{j\ne i} k(l_{ij})\, s_{ij}\, \ee^{\,e}_{ij} \;=\; \mathbf{0},
\qquad i = 1,\ldots,N,
\label{eq:equilibrium-s}
\end{equation}
so the forces on each node cancel, which does not require the individual $s_{ij}$
to vanish. Substituting \eqref{eq:sij} into \eqref{eq:V} evaluated at $\re$ gives
the minimum energy
\begin{equation}
\boxed{\;
V(\re; \LL) \;=\; \tfrac12 \sum_{i<j} k(l_{ij})\, s_{ij}^{2},
\;}
\label{eq:Vmin}
\end{equation}
which depends on the strains only through their squares, and so is unchanged by the
sign convention adopted in \eqref{eq:sij}.

\subsection{Frustration}
\label{sec:frustration}

We call a network \textbf{frustrated} when $V(\re;\LL) > 0$: no conformation
relaxes all of its active springs at once, and the network sits at a compromise in
which the forces cancel node by node while individual springs remain strained. By
\eqref{eq:Vmin} this is the case whenever $s_{ij} \ne 0$ for at least one active
spring. A network that is not frustrated has $s_{ij} = 0$ for every active spring,
and then $V(\re;\LL) = 0$.

Whether a given $\LL$ is frustrated is settled by solving \eqref{eq:equilibrium}
for $\re$ and evaluating \eqref{eq:Vmin}. One case can be read off without
solving anything. Choose a structure $\rr^{\ast}$ and set
$l_{ij} = d_{ij}(\rr^{\ast})$ for every spring. Then every strain vanishes at
$\rr^{\ast}$, so \eqref{eq:equilibrium-s} is satisfied term by term,
$\rr^{\ast}$ is an equilibrium conformation, and $V(\rr^{\ast};\LL) = 0$: the
network is unfrustrated by construction. This is what a standard elastic network
model does, and it is a statement about how $\LL$ was chosen, not about
\eqref{eq:equilibrium}, which takes $\LL$ as given.

Nothing obliges $\LL$ to have been chosen that way, and parameter sets not of this
form will in general be frustrated. How the $l_{ij}$ ought to be chosen is a
separate question from the one \eqref{eq:equilibrium} answers, and is not taken up
here; Section~\ref{sec:open} returns to it. What matters below is that both cases
are allowed: nothing in Sections \ref{sec:distribution} to \ref{sec:changes}
assumes an unfrustrated network. As will be seen in Section~\ref{sec:mutation}, a
mutation applied to an unfrustrated network produces a frustrated one.

\section{The conformational distribution}
\label{sec:distribution}

\subsection{The Boltzmann distribution}

At temperature $T$ the conformations are distributed according to
\begin{equation}
\rho(\rr) \;=\; \frac{e^{-\beta V(\rr;\LL)}}{Z},
\qquad
Z \;=\; \int e^{-\beta V(\rr;\LL)}\, d\rr,
\qquad
\beta = \frac{1}{k_B T},
\label{eq:boltzmann}
\end{equation}
where the integral runs over $\mathbb{R}^{3N}$. This is a configurational
integral: the momenta have been left out. They contribute a factor that depends on
the masses and on $\beta$ but not on $\LL$, so it cancels from every difference
between networks computed in Section~\ref{sec:changes}, which is all we use.

\subsection{The Hessian}

The integral in \eqref{eq:boltzmann} cannot be evaluated in closed form, because
$V$ is not quadratic in $\rr$: the distances \eqref{eq:dij} are nonlinear
functions of the coordinates. We therefore expand $V$ about an equilibrium
conformation. Writing $\mathbf{u} = \rr - \re$,
\begin{equation}
V(\rr;\LL) \;=\; V(\re;\LL)
\;+\; \frac{\partial V}{\partial \rr}\bigg|_{\re}\!\!\cdot \mathbf{u}
\;+\; \tfrac12\, \mathbf{u}^{\mathsf T} \KK\, \mathbf{u} \;+\; O(\mathbf{u}^3),
\qquad
\KK \;=\; \frac{\partial^2 V}{\partial \rr^2}\bigg|_{\re}.
\label{eq:expansion}
\end{equation}
The linear term vanishes by \eqref{eq:equilibrium}.

To obtain $\KK$ we need the second derivatives of the distances. Differentiating
$\ee_{ij} = \rr_{ij}/d_{ij}$ with respect to $\rr_j$, using the product rule with
$\partial \rr_{ij}/\partial\rr_j = \mathbf{I}$ and
$\partial d_{ij}^{-1}/\partial \rr_j = -d_{ij}^{-2}\,\ee_{ij}^{\mathsf T}$,
\begin{equation}
\frac{\partial \ee_{ij}}{\partial \rr_j}
\;=\; \frac{\mathbf{I}}{d_{ij}} - \frac{\rr_{ij}\ee_{ij}^{\mathsf T}}{d_{ij}^{2}}
\;=\; \frac{\mathbf{I} - \ee_{ij}\ee_{ij}^{\mathsf T}}{d_{ij}}.
\label{eq:d2d}
\end{equation}
Now differentiate the gradient \eqref{eq:grad} with respect to $\rr_j$ for $j \ne i$.
Only the single term of the sum with that index depends on $\rr_j$, and it is a
product of $(d_{ij} - l_{ij})$ and $\ee_{ij}$, so
\begin{equation}
\KK_{ij}
\;=\; -\,k(l_{ij})\left[
\frac{\partial\, (d_{ij} - l_{ij})}{\partial \rr_j}\,\ee_{ij}^{\mathsf T}
\;+\; (d_{ij} - l_{ij})\, \frac{\partial \ee_{ij}}{\partial \rr_j}
\right].
\end{equation}
Substituting $\partial d_{ij}/\partial\rr_j = \ee_{ij}$ and \eqref{eq:d2d}, and
then $d_{ij} - l_{ij} = s_{ij}$ from \eqref{eq:sij},
\begin{equation}
\KK_{ij}
\;=\; -\,k(l_{ij})\left[ \ee_{ij}\ee_{ij}^{\mathsf T}
+ s_{ij}\,\frac{\mathbf{I} - \ee_{ij}\ee_{ij}^{\mathsf T}}{d_{ij}} \right],
\end{equation}
that is, defining $g_{ij} \equiv s_{ij}/d_{ij}(\re)$,
\begin{equation}
\boxed{\;
\KK_{ij} \;=\; -\,k(l_{ij})\Big[\,\ee_{ij}\ee_{ij}^{\mathsf T}
- g_{ij}\big(\ee_{ij}\ee_{ij}^{\mathsf T} - \mathbf{I}\big)\Big],
\qquad i \ne j.
\;}
\label{eq:hessian}
\end{equation}
The diagonal blocks follow from translational invariance. Displacing every node by
the same vector $\mathbf{a}$ leaves $V$ unchanged, so $\KK$ annihilates that
displacement, and summing the blocks in a row gives
$\sum_{j} \KK_{ij}\mathbf{a} = \mathbf{0}$ for every $\mathbf{a}$, whence
\begin{equation}
\KK_{ii} \;=\; -\sum_{j \ne i} \KK_{ij}.
\label{eq:Kii}
\end{equation}
Inactive springs have $k(l_{ij}) = 0$ and contribute nothing, so $\KK$ is
$3N \times 3N$ however many springs carry parameters.

\subsection{Longitudinal and transverse stiffness}

Equation \eqref{eq:hessian} contains two projectors. Writing
$\mathbf{P}_\perp = \mathbf{I} - \ee_{ij}\ee_{ij}^{\mathsf T}$ for the projector
onto the plane perpendicular to the spring, it may be rearranged as
\begin{equation}
\KK_{ij} \;=\; -\,k(l_{ij})\Big[\ee_{ij}\ee_{ij}^{\mathsf T}
+ g_{ij}\,\mathbf{P}_\perp\Big],
\label{eq:split}
\end{equation}
so a spring carries two stiffnesses rather than one: a longitudinal $k(l_{ij})$
resisting change of the distance, and a transverse $k(l_{ij})g_{ij}$, doubly
degenerate, resisting rotation of the spring at fixed distance.

The transverse stiffness is proportional to $g_{ij} = s_{ij}/d_{ij}(\re)$, and so
vanishes for a relaxed spring and changes sign with the strain. A stretched spring
($s_{ij}>0$) resists transverse displacement, as a taut string does. A compressed
one ($s_{ij}<0$) has negative transverse stiffness: displacing the two nodes
sideways increases their separation, which a compressed spring favours. Note that
it is the transverse stiffness that is negative, not $k(l_{ij})$, which is
non-negative by \eqref{eq:kl}.

Since $g_{ij} = 0$ for every spring of an unfrustrated network, the transverse term
is absent there, and \eqref{eq:hessian} reduces to the rank-one form of a standard
elastic network model. Frustration enters the conformational distribution through
$g_{ij}$ and through nothing else.

\subsection{The Gaussian approximation}

Truncating \eqref{eq:expansion} after the quadratic term gives
\begin{equation}
V(\rr;\LL) \;\simeq\; V(\re;\LL) + \tfrac12 \mathbf{u}^{\mathsf T}\KK\mathbf{u},
\label{eq:quadratic}
\end{equation}
and hence, from \eqref{eq:boltzmann}, a Gaussian distribution of conformations
centred on $\re$,
\begin{equation}
\rho(\rr) \;\simeq\; \frac{1}{Z}\, e^{-\beta V(\re;\LL)}\,
\exp\!\big(-\tfrac{\beta}{2}\,\mathbf{u}^{\mathsf T}\KK\,\mathbf{u}\big).
\label{eq:gaussian}
\end{equation}

$\KK$ is symmetric, and positive semi-definite because $\re$ is a minimum. Let its
eigenvalues be $\lambda_n$ with eigenvectors $\mathbf{q}_n$. At least six
eigenvalues vanish, corresponding to the three translations and three rotations
along which $V$ does not vary.

Whether there are exactly six is a property of the network, not a given. A
displacement $\mathbf{u}$ is a null direction of $\KK$ when it changes no active
distance to first order, that is when $\ee_{ij}\cdot(\mathbf{u}_i - \mathbf{u}_j) = 0$
for every active spring. A network for which the only such displacements are the six
rigid motions is \emph{infinitesimally rigid}, and then $\mathrm{rank}\,\KK = 3N-6$.
Connectivity does not suffice: a node joined to the rest by a single active spring
can swing about it, and one joined by two can rotate about the axis through them,
in both cases without changing any active distance. Each such freedom is a further
null direction. In what follows the networks are assumed infinitesimally rigid,
and we write $\sum'$ and $\prod'$ for sums and products over the $3N-6$ nonzero
eigenvalues.

The distribution \eqref{eq:gaussian} is flat along the six null directions, which
is the rigid-body freedom of the whole network, and we factor it out: displacements
are taken in the complement, where the covariance is
\begin{equation}
\bm{\Sigma} \;=\; \beta^{-1}\KK^{+},
\qquad
\KK^{+} \;=\; \sum_n{}' \lambda_n^{-1}\, \mathbf{q}_n \mathbf{q}_n^{\mathsf T},
\label{eq:covariance}
\end{equation}
with $\KK^{+}$ the pseudo-inverse obtained by inverting the nonzero eigenvalues
and annihilating the null directions.

\section{Thermodynamic properties}
\label{sec:thermo}

\subsection{The partition function}

In the quadratic approximation \eqref{eq:quadratic} the integral
\eqref{eq:boltzmann} factorises in the eigenbasis of $\KK$. Writing
$\mathbf{u} = \sum_n y_n \mathbf{q}_n$, the quadratic form is
$\sum_n \lambda_n y_n^2$, and the integral over each of the $3N-6$ non-null
directions is a standard Gaussian,
\begin{equation}
\int_{-\infty}^{\infty} e^{-\beta\lambda_n y_n^2/2}\, dy_n
\;=\; \left(\frac{2\pi}{\beta\lambda_n}\right)^{1/2}.
\end{equation}
Hence
\begin{equation}
\boxed{\;
Z \;=\; e^{-\beta V(\re;\LL)} \prod_n{}' \left(\frac{2\pi}{\beta\lambda_n}\right)^{1/2},
\qquad
\ln Z \;=\; -\beta V(\re;\LL) + \tfrac12 \sum_n{}' \ln\frac{2\pi}{\beta\lambda_n}.
\;}
\label{eq:Z}
\end{equation}

\subsection{Free energy, internal energy, entropy}

The Helmholtz free energy is
\begin{equation}
A \;=\; -\beta^{-1}\ln Z
\;=\; V(\re;\LL) \;-\; \frac{1}{2\beta}\sum_n{}' \ln\frac{2\pi}{\beta\lambda_n}.
\label{eq:A}
\end{equation}
The internal energy follows by differentiating \eqref{eq:Z}, using
$\partial\ln\lambda_n/\partial\beta = 0$,
\begin{equation}
U \;=\; -\frac{\partial \ln Z}{\partial \beta}
\;=\; V(\re;\LL) \;+\; \frac{3N-6}{2\beta},
\label{eq:U}
\end{equation}
which is the minimum energy plus $\tfrac12 k_B T$ per non-rigid degree of freedom,
as equipartition requires. The entropy is $S = (U-A)/T$, so
\begin{equation}
T S \;=\; U - A \;=\; \frac{1}{2\beta}\sum_n{}'
\left[\ln\frac{2\pi}{\beta\lambda_n} + 1\right].
\label{eq:S}
\end{equation}
Soft directions, those of small $\lambda_n$, contribute most: the network gains
entropy by being easy to deform.

Two comments on \eqref{eq:A} and \eqref{eq:S}. The logarithm takes a dimensional
argument, so the numerical values of $A$ and $TS$ depend on the units chosen for
length, and $A$ may fall below $V(\re;\LL)$. Neither is meaningful alone. What is
meaningful is a difference between two networks at the same temperature, where the
units cancel; those are the quantities of Section~\ref{sec:changes}. The same
applies to the momentum factor dropped in \eqref{eq:boltzmann}, and to the
rigid-body volume factored out of \eqref{eq:covariance}: both are common to the
two networks and cancel.

\jc{Open question. Equations \eqref{eq:A} and \eqref{eq:S} rest on the quadratic
truncation \eqref{eq:quadratic}. How good it is for a frustrated network, where the
transverse term makes $\KK$ differ from the standard form, has not been examined
here.}

\section{Reference and mutant}
\label{sec:mutation}

\subsection{The reference network}

Everything above applies to any network. We now fix one, the \emph{reference},
with parameters $\LL^{\text{ref}} = \{l^{\text{ref}}_{ij}\}$ and potential
\begin{equation}
V_{\text{ref}}(\rr) \;\equiv\; V(\rr; \LL^{\text{ref}})
\;=\; \tfrac12 \sum_{i<j} k\big(l^{\text{ref}}_{ij}\big)
\Big(d_{ij}(\rr) - l^{\text{ref}}_{ij}\Big)^{2}.
\label{eq:Vref}
\end{equation}
Write $k^{\text{ref}}_{ij} = k(l^{\text{ref}}_{ij})$ for its force constants,
$\re_{\text{ref}}$ for an equilibrium conformation obtained from
\eqref{eq:equilibrium}, and
\begin{equation}
s_{ij} \;=\; d_{ij}(\re_{\text{ref}}) - l^{\text{ref}}_{ij}
\label{eq:sref}
\end{equation}
for its strains. Its Hessian $\KK^{\text{ref}}$, spectrum
$\{\lambda^{\text{ref}}_n\}$, and thermodynamic properties $A_{\text{ref}}$,
$U_{\text{ref}}$, $S_{\text{ref}}$ follow from Sections \ref{sec:distribution} and
\ref{sec:thermo}. The reference may be frustrated or not; nothing below assumes
either.

\subsection{Mutation}

A mutation at site $m$ perturbs the equilibrium lengths of the springs that
connect $m$ to the rest of the network, and leaves all others unchanged:
\begin{equation}
l^{\text{mut}}_{ij} \;=\; l^{\text{ref}}_{ij} + \delta l_{ij},
\qquad
\delta l_{ij} = 0 \ \text{ unless } i = m \text{ or } j = m.
\label{eq:mutation}
\end{equation}
The perturbations $\delta l_{ij}$ are drawn from a distribution with zero mean and
variance $\sigma^2$, independently for each affected spring; $\sigma$ sets the size
of the mutation.

The mutant potential is \eqref{eq:V} evaluated with these parameters,
\begin{equation}
V_{\text{mut}}(\rr) \;\equiv\; V(\rr; \LL^{\text{mut}})
\;=\; \tfrac12 \sum_{i<j} k\big(l^{\text{ref}}_{ij} + \delta l_{ij}\big)
\Big(d_{ij}(\rr) - l^{\text{ref}}_{ij} - \delta l_{ij}\Big)^{2}.
\label{eq:Vmut}
\end{equation}
Comparing \eqref{eq:Vref} and \eqref{eq:Vmut}: the mutation changes both the
equilibrium lengths and, through \eqref{eq:kl}, the force constants. Define
\begin{equation}
\delta k_{ij} \;\equiv\; k\big(l^{\text{ref}}_{ij} + \delta l_{ij}\big)
- k\big(l^{\text{ref}}_{ij}\big),
\qquad\text{so}\qquad
k^{\text{mut}}_{ij} = k^{\text{ref}}_{ij} + \delta k_{ij}.
\label{eq:dk}
\end{equation}
This is an exact difference of two values of $k$, not a differential: no
assumption is made that $\delta l_{ij}$ is small or that $k$ is differentiable.
A spring whose equilibrium length crosses from a region where $k>0$ to one where
$k=0$, or the reverse, has $\delta k_{ij}$ of the size of $k$ itself. Contacts are
created and destroyed here, in the parameters, and not by any change in
conformation.

This is why every spring carries an equilibrium length, the inactive ones
included. An inactive spring with no $l_{ij}$ could never acquire one inside the
range where $k > 0$, so a mutation could only destroy contacts and never form
them; carrying the parameter regardless lets \eqref{eq:mutation} move a spring in
either direction across that range.

The mutant has its own equilibrium conformation $\re_{\text{mut}}$, obtained from
\eqref{eq:equilibrium} with the mutant parameters, its own strains
$s^{\text{mut}}_{ij} = d_{ij}(\re_{\text{mut}}) - l^{\text{mut}}_{ij}$, its own
Hessian $\KK^{\text{mut}}$ and spectrum $\{\lambda^{\text{mut}}_n\}$, and its own
$A_{\text{mut}}$, $U_{\text{mut}}$, $S_{\text{mut}}$. Note that $\re_{\text{mut}}$
is obtained by minimising \eqref{eq:Vmut} itself, not by displacing
$\re_{\text{ref}}$ along a response to the perturbation.

Even if the reference is unfrustrated, the mutant will in general be frustrated:
the perturbed lengths $l^{\text{mut}}_{ij}$ need not be simultaneously realisable
by any conformation, so the mutant settles at a compromise with $s^{\text{mut}}_{ij}
\ne 0$.

\section{Mutational changes}
\label{sec:changes}

\subsection{Three energy differences}

Two networks and two conformations give three energy differences worth naming. The
first is the cost of changing the parameters at fixed conformation,
\begin{equation}
\Delta V_{\text{stress}} \;=\;
V_{\text{mut}}(\re_{\text{ref}}) - V_{\text{ref}}(\re_{\text{ref}}),
\label{eq:dvstress}
\end{equation}
the second the energy released as the mutant relaxes to its own equilibrium,
\begin{equation}
\Delta V_{\text{relax}} \;=\;
V_{\text{mut}}(\re_{\text{mut}}) - V_{\text{mut}}(\re_{\text{ref}}),
\label{eq:dvrelax}
\end{equation}
and their sum is the difference between the two minima,
\begin{equation}
\Delta V_{\min} \;=\;
V_{\text{mut}}(\re_{\text{mut}}) - V_{\text{ref}}(\re_{\text{ref}})
\;=\; \Delta V_{\text{stress}} + \Delta V_{\text{relax}}.
\label{eq:dvmin}
\end{equation}
Since $\re_{\text{mut}}$ minimises $V_{\text{mut}}$, we have
$V_{\text{mut}}(\re_{\text{mut}}) \le V_{\text{mut}}(\rr)$ for every $\rr$, and in
particular for $\rr = \re_{\text{ref}}$, so
\begin{equation}
\Delta V_{\text{relax}} \;\le\; 0 .
\end{equation}

\subsection{Decomposition of the stress energy}
\label{sec:decomposition}

$\Delta V_{\text{stress}}$ can be written out in closed form. We do so in steps.

\paragraph{Step 1: separate the springs.}
Both \eqref{eq:Vref} and \eqref{eq:Vmut} are sums over the same index set, all
$\tfrac12 N(N-1)$ springs, so their difference is a sum of one term per spring,
\begin{equation}
\Delta V_{\text{stress}} \;=\; \sum_{i<j} \left[
\tfrac12 k^{\text{mut}}_{ij}\big(d^e_{ij} - l^{\text{mut}}_{ij}\big)^2
- \tfrac12 k^{\text{ref}}_{ij}\big(d^e_{ij} - l^{\text{ref}}_{ij}\big)^2
\right],
\qquad d^e_{ij} \equiv d_{ij}(\re_{\text{ref}}).
\end{equation}
This is where carrying a parameter for every spring, active or not, does its work:
no spring appears in one sum and not the other, and a contact created or destroyed
by the mutation is carried by $\delta k_{ij}$ rather than by a change in the range
of summation.

\paragraph{Step 2: eliminate the distance.}
By the definition of the reference strain \eqref{eq:sref},
$d^e_{ij} = l^{\text{ref}}_{ij} + s_{ij}$, so the two brackets become
\begin{equation}
d^e_{ij} - l^{\text{ref}}_{ij} \;=\; s_{ij},
\qquad
d^e_{ij} - l^{\text{mut}}_{ij} \;=\; s_{ij} - \delta l_{ij}.
\end{equation}
The conformation has now disappeared from the problem: what remains is expressed
entirely in the reference strains and the perturbation.

\paragraph{Step 3: substitute the force constants.}
Using \eqref{eq:dk}, the term for one spring is
\begin{equation}
\tfrac12 \big(k^{\text{ref}}_{ij} + \delta k_{ij}\big)
\big(s_{ij} - \delta l_{ij}\big)^2
- \tfrac12 k^{\text{ref}}_{ij}\, s_{ij}^2 .
\end{equation}

\paragraph{Step 4: split off the change in stiffness.}
Grouping the two terms in $k^{\text{ref}}_{ij}$ and keeping the $\delta k_{ij}$
term apart,
\begin{equation}
\;=\; \tfrac12 k^{\text{ref}}_{ij}\Big[\big(s_{ij} - \delta l_{ij}\big)^2 - s_{ij}^2\Big]
\;+\; \tfrac12\, \delta k_{ij} \big(s_{ij} - \delta l_{ij}\big)^2 .
\end{equation}
The second piece is already in final form. No expansion has been made in obtaining
it, and none is possible: $\delta k_{ij}$ is not assumed small.

\paragraph{Step 5: expand the square.}
The bracket is a difference of two squares,
\begin{equation}
\big(s_{ij} - \delta l_{ij}\big)^2 - s_{ij}^2
\;=\; -\,2\, s_{ij}\, \delta l_{ij} + \delta l_{ij}^2 ,
\end{equation}
which is a finite algebraic identity, not a truncated expansion: it is exact for
any $\delta l_{ij}$.

\paragraph{Step 6: collect.}
Summing over springs,
\begin{equation}
\boxed{\;
\Delta V_{\text{stress}} \;=\;
\underbrace{\tfrac12 \sum_{i<j} k^{\text{ref}}_{ij}\, \delta l_{ij}^{2}}_{\text{(a)}}
\;-\; \underbrace{\sum_{i<j} k^{\text{ref}}_{ij}\, s_{ij}\, \delta l_{ij}}_{\text{(b)}}
\;+\; \underbrace{\tfrac12 \sum_{i<j} \delta k_{ij}
\big(s_{ij} - \delta l_{ij}\big)^{2}}_{\text{(c)}} .
\;}
\label{eq:decomposition}
\end{equation}
This holds for any $k$, any mutation size, and any reference network, frustrated or
not.

\subsection{The three terms}

Term (a) comes from the $\delta l_{ij}^2$ of the expanded square. It is
non-negative, quadratic in the size of the mutation, and independent of the state
of the reference.

Term (b) is the cross term $-2 s_{ij}\delta l_{ij}$ of the same square. It is
proportional to the reference strain and is the only term of the three that can be
negative. It contributes $-k^{\text{ref}}_{ij} s_{ij}\,\delta l_{ij}$ per spring,
which is negative when $s_{ij}$ and $\delta l_{ij}$ share a sign: lengthening the
equilibrium length of a stretched spring ($s_{ij} > 0$, $\delta l_{ij} > 0$), or
shortening that of a compressed one, moves the equilibrium length towards the
distance the spring actually spans and so releases energy. The opposite pairing
costs energy. Term (b) vanishes when $s_{ij}\,\delta l_{ij} = 0$ for every spring,
which is to say when the reference is unstrained on the springs the mutation
touches. In particular it vanishes for any mutation of an unfrustrated reference,
and it does not vanish in general.

Term (c) is the whole of the contribution from the change in force constants. It
vanishes when $\delta k_{ij}(s_{ij} - \delta l_{ij})^2 = 0$ for every spring, in
particular when $k$ is constant over the range the mutation explores. It is the
term that a model with a fixed contact map does not have.

Terms (b) and (c) are what distinguish \eqref{eq:decomposition} from the energy
cost of mutating an unfrustrated network with fixed force constants, which is (a)
alone. Neither is a correction of higher order in the mutation size: both are
exact contributions of the same standing as (a), which happen to vanish in that
case.

\subsection{Differences and energies}

$\Delta V_{\text{stress}}$ is a difference of two potentials, and is not the same
as the absolute energy $V_{\text{mut}}(\re_{\text{ref}})$ of the mutant network at
the reference conformation. By \eqref{eq:dvstress} the two differ by
$V_{\text{ref}}(\re_{\text{ref}})$, the strain energy the reference already
carries, which by \eqref{eq:Vmin} is zero for an unfrustrated reference and
positive otherwise.

\subsection{Changes in the thermodynamic properties}

Applying Section~\ref{sec:thermo} to both networks and subtracting, the internal
energy difference follows from \eqref{eq:U},
\begin{equation}
\Delta U \;=\; U_{\text{mut}} - U_{\text{ref}}
\;=\; V_{\text{mut}}(\re_{\text{mut}}) - V_{\text{ref}}(\re_{\text{ref}})
\;=\; \Delta V_{\min},
\end{equation}
the equipartition terms cancelling. This step, and the entropy difference below,
require the two networks to have the same number of non-null directions. A point
mutation changes the parameters of one site and neither adds nor removes a node,
so both Hessians are $3N\times3N$; and if both networks are infinitesimally rigid
in the sense of Section~\ref{sec:distribution}, both have $3N-6$ nonzero
eigenvalues. For the entropy, write \eqref{eq:S} as
\begin{equation}
TS \;=\; \frac{1}{2\beta}\left[ (3N-6)\left(\ln\frac{2\pi}{\beta} + 1\right)
- \sum_n{}' \ln \lambda_n \right],
\end{equation}
in which only the last term depends on the network. Both networks have the same
number of non-null directions, so on subtracting, the constant term cancels and
\begin{equation}
\boxed{\;
T \Delta S \;=\; \frac{1}{2\beta}\left[
\sum_n{}' \ln \lambda^{\text{ref}}_n - \sum_n{}' \ln \lambda^{\text{mut}}_n \right]
\;=\; \frac{1}{2\beta} \ln
\frac{\prod_n' \lambda^{\text{ref}}_n}{\prod_n' \lambda^{\text{mut}}_n} .
\;}
\label{eq:dS}
\end{equation}
The two spectra enter only through these products. No correspondence between the
modes of the reference and those of the mutant is used, and none is available:
the eigenvectors of $\KK^{\text{ref}}$ and $\KK^{\text{mut}}$ are different sets of
directions, and the eigenvalues are not paired by any property of the model.

The free energy difference is then
\begin{equation}
\Delta A \;=\; \Delta U - T\Delta S
\;=\; \Delta V_{\min} - \frac{1}{2\beta} \ln
\frac{\prod_n' \lambda^{\text{ref}}_n}{\prod_n' \lambda^{\text{mut}}_n}.
\label{eq:dA}
\end{equation}
The dimensional ambiguity noted after \eqref{eq:S} has cancelled: \eqref{eq:dS} and
\eqref{eq:dA} depend only on the ratio of the two products. In the protein
literature the quantity \eqref{eq:dA} is what is meant by $\Delta\Delta G$ for the
folded state.

The products $\prod_n' \lambda_n$ are the pseudo-determinants of the two Hessians,
so \eqref{eq:dS} may be computed without diagonalising either matrix.

Equations \eqref{eq:dvmin} and \eqref{eq:dS} separate the two ways a mutation
changes the free energy. It shifts the minimum of the potential, which
\eqref{eq:decomposition} decomposes; and it changes the spectrum of $\KK$, which
by \eqref{eq:hessian} depends on the strains through $g_{ij}$ and on the contact
map through $k(l_{ij})$. A mutation that leaves the minimum energy untouched can
still change the free energy, through the spectrum alone.

\section{The step function}

The results above hold for any $k$. The common choice is a step at a cutoff
distance,
\begin{equation}
k(l) \;=\; k_0\, \mathds{1}\big[\, l \le d_{\max} \,\big],
\label{eq:step}
\end{equation}
under which a spring is either active with strength $k_0$ or inactive. Two
consequences follow from the results above.

First, $V$ is discontinuous in the parameters. If a mutation carries
$l^{\text{ref}}_{ij}$ across $d_{\max}$ then $\delta k_{ij} = \pm k_0$, and by term
(c) of \eqref{eq:decomposition} the energy changes by
$\pm\tfrac12 k_0 (s_{ij} - \delta l_{ij})^2$ however small $\delta l_{ij}$ is.

Second, and for the same reason, term (c) cannot be treated as small when
$\delta l_{ij}$ is small. This is the sense in which the contact map resists a
perturbative treatment: $\delta k_{ij}$ is not controlled by $\delta l_{ij}$.

A continuous $k$, for instance a sigmoid of width $w$ centred on $d_{\max}$, makes
$V$ continuous in $\LL$ and $\delta k_{ij}$ small whenever $\delta l_{ij}$ is
small. The cost is that no spring is exactly inactive, so the contact map becomes a
threshold imposed on a continuum rather than a property of the network.

\section{Open questions and work to come}
\label{sec:open}

The following are not settled here.

\begin{itemize}
\item The quadratic truncation \eqref{eq:quadratic} is the only approximation in
the document. Its accuracy for a frustrated network has not been examined.
\item The magnitudes of (a), (b) and (c) in \eqref{eq:decomposition} for a real
protein are not known. The decomposition says when each vanishes, not how large
each is.
\item How large $|s_{ij}|$ may become before a compressed spring ($s_{ij} < 0$),
with its negative transverse stiffness \eqref{eq:split}, ceases to represent a
residue contact.
\item A mutation can cost the network its infinitesimal rigidity. The perturbation
\eqref{eq:mutation} acts on every spring of the mutated site, and if enough of
their equilibrium lengths leave the range where $k>0$, the site is left attached
by one or two springs and acquires a rotational freedom that changes no active
distance. The mutant Hessian then has more than six null directions, the counts
in Section~\ref{sec:changes} no longer match, and neither $\Delta U =
\Delta V_{\min}$ nor \eqref{eq:dS} holds as written: an uncancelled term in
$\ln(2\pi/\beta)$ survives, so $\Delta A$ is no longer independent of the units of
length. How often this happens at a given mutation size and cutoff is an empirical
question, and what the free energy difference should mean when it does is not
settled here.
\item How the parameters $\LL$ should be chosen. Section~\ref{sec:frustration}
gives one recipe, $l_{ij} = d_{ij}(\rr^{\ast})$ for a chosen structure, which
yields an unfrustrated network; but nothing selects it, and the question of
whether a frustrated network is the better model of a native structure --- and if
so, which one --- is open. This is a question about the parameters, and is not
one \eqref{eq:equilibrium} addresses: that equation takes $\LL$ as given and
solves for the conformation.
\end{itemize}

To be added: the distribution of structures generated by repeated mutation, and
the entropic contribution to $\Delta\Delta G$ along a substitution process.

\end{document}
